%Using Preamble.tex of paulgp/draft-tex
%https://github.com/paulgp/draft-tex

\usepackage[top=0.8in, bottom=0.78in, left=0.87in, right=0.87in]{geometry}
\usepackage{setspace}
\usepackage[T1]{fontenc}
\usepackage{times}
\usepackage{booktabs}
\usepackage{rotating}
\usepackage{graphicx}
\usepackage[section]{placeins} %Placeins.sty keeps floats `in their place', preventing them from floating past a "\FloatBarrier" command into another section.  To use it, declare "\usepackage{placeins}" and insert "\FloatBarrier" at places that floats should not move past, perhaps at every "\section".  
\usepackage[large, bf]{caption}
\usepackage[FIGTOPCAP]{subfigure}
\usepackage{pdfpages}
\usepackage{palatino}
\usepackage{pdflscape}
\usepackage{textcomp}
\usepackage{longtable}
\usepackage{nicefrac}
\usepackage{adjustbox}	% to adjust the size of objects to fit into a page
\usepackage[hyphens]{url}


\usepackage{natbib}
\bibpunct{(}{)}{;}{a}{,}{,}

\def\citeapos#1{\citeauthor{#1}'s (\citeyear{#1})}

\doublespacing

\usepackage{amsmath, amsfonts, amssymb, amsthm}

\usepackage{mathpazo} %Use Palotino fonts
\parskip 0ex  %Vertical distance between paragraphs, in "ex"s
\parindent 20pt


\usepackage[pdftex]{hyperref}
\hypersetup{colorlinks, citecolor=blue, filecolor=blue, linkcolor=blue, urlcolor=blue}



%\usepackage{subcaption} % for figures


\newtheorem{theorem}{Theorem}
\newtheorem{lemma}{Lemma}
\newtheorem{proposition}{Proposition}
\newtheorem{corollary}{Corollary}
\newtheorem{prediction}{Prediction}
\newtheorem{case}{Special Case}

\newenvironment{proofAlt}[1][Proof]{\begin{trivlist}
\item[\hskip \labelsep {\bfseries #1}]}{\end{trivlist}}
\newenvironment{definition}[1][Definition]{\begin{trivlist}
\item[\hskip \labelsep {\bfseries #1}]}{\end{trivlist}}
\newenvironment{example}[1][Example]{\begin{trivlist}
\item[\hskip \labelsep {\bfseries #1}]}{\end{trivlist}}

\newenvironment{remark}[1][Remark]{\begin{trivlist}
\item[\hskip \labelsep {\bfseries #1}]}{\end{trivlist}}
\def\urltilda{\kern -.15em\lower .7ex\hbox{\~{}}\kern .04em}

\renewcommand{\thesubfigure}{(\Alph{subfigure})}


%%%%%%%%%%%%%%%%%%%%%%%%%%%%%%%%%%
% SPACING
%%%%%%%%%%%%%%%%%%%%%%%%%%%%%%%%%%

\usepackage{titlesec}

\titlespacing*{\section}{0pt}{1.5ex plus 1ex minus .2ex}{0.8ex plus .2ex}
\titlespacing*{\subsection}{0pt}{1.2ex plus 1ex minus .2ex}{0.8ex plus .2ex}
\usepackage{appendix}
% Decimal align 
\usepackage{dcolumn}
\newcolumntype{d}[0]{D{.}{.}{5}}

% Spaces in file path
\usepackage{graphicx}
\usepackage[space]{grffile}

%%%%%%%%%%%%%%%%%%%%%%%%%%%%%%%%%%
% Tables
\usepackage{threeparttable}

%%%%%%%%%%%%%%%%%%%%%%%%%%%%%%%%%%
% Packages for Tables generated by R ----
\usepackage{tabularray}
\usepackage{booktabs} 
\usepackage{longtable}
\usepackage{array}
\usepackage{multirow}
\usepackage{wrapfig} % if necessary
\usepackage{float}
\usepackage{colortbl} % if necessary
\usepackage{pdflscape}
\usepackage{tabu}
\usepackage{threeparttable}
\usepackage{threeparttablex}
\usepackage[normalem]{ulem}
\usepackage{makecell}
\usepackage{xcolor} % if necessary
\usepackage{siunitx}
\usepackage{comment}
% -----------------------------------------

%%%%%%%%%%%%%%%%%%%%%%%%%%%%%%%%%%%%%%%%%%%%%%%%%%%%%%%%%%%%%

%%%%%%%%%%%%%%%%%%%%%%%%%%%%%%%%
%For long notes
%https://www.hargaden.com/enda/long-or-multiple-line-notes-in-esttab-with-automatic-wrapping/
\newcommand{\tabnotes}[2]{\bottomrule \multicolumn{#1}{@{}p{0.70\linewidth}@{}}{\footnotesize #2 }\end{tabular}\end{table}} 
%%%%%%%%%%%%%%%%%%%%%%%%%%%%%%%%
